% ================= RESUMEN EJECUTIVO =================
\section*{Resumen ejecutivo}
\addcontentsline{toc}{section}{Resumen ejecutivo}

La inteligencia artificial jurídica en español ---y en particular para el derecho mexicano--- se encuentra estructuralmente desatendida. Los modelos de lenguaje legales dominantes se entrenan sobre corpus angloamericanos, alucinan citas cuando se les pregunta sobre legislación mexicana y no pueden desplegarse localmente para despachos obligados al secreto profesional. La evidencia empírica disponible es contundente: los modelos generales de propósito amplio alucinan entre el 58\% y el 88\% de las respuestas a consultas jurídicas verificables,\cite{dahl2024} e incluso las herramientas jurídicas comerciales con recuperación aumentada (RAG) ---Lexis+ AI y Westlaw AI--- alucinan entre el 17\% y el 33\% de las veces.\cite{magesh2025} Ninguno de estos sistemas está anclado en fuentes primarias mexicanas.

Este whitepaper presenta \textbf{AI Justicia}, un esfuerzo integral para construir el primer sistema de IA legal soberano de México, y evalúa de manera exhaustiva su necesidad, factibilidad técnica y viabilidad de negocio. El proyecto descansa sobre tres pilares: (i) la construcción de un \textbf{corpus vivo de más de 10B tokens de fuentes primarias exclusivamente mexicanas} ---legislación federal y estatal, gacetas oficiales, jurisprudencia federal y local, y doctrina universitaria---, cuya base normativa de reutilización es sólida (el art.~14, fracc.~VII de la LFDA excluye los textos legislativos, administrativos y judiciales de la protección autoral\cite{lfda}); (ii) un \textbf{harness de generación anclada en recuperación} donde toda afirmación jurídica debe verificarse contra pasajes recuperados, con abstención honesta cuando la evidencia es insuficiente; y (iii) una \textbf{estrategia de entrenamiento soberano} para \textbf{Tlamatini} ---el modelo jurídico mexicano del proyecto, del náhuatl \emph{tlamatini}, el que sabe'': los sabios-consejeros del México central prehispánico--- que sintetiza las lecciones de SaulLM-7B (continued pretraining de dominio con replay)\cite{colombo2024} y de la generación industrial de modelos legales (CPT fuerte más fusión lineal), extendida con \emph{harness training}: enseñar el formato de respuesta con citas ancladas durante el propio preentrenamiento.

\begin{tesis}
\textbf{Tesis central.} Para IA jurídica específica de jurisdicción, la combinación de \emph{harness en los pesos} más \emph{recuperación en inferencia} supera a cualquiera de las dos estrategias por separado. El CPT de dominio produce la \emph{forma} correcta de la respuesta jurídica; la recuperación produce los \emph{hechos} correctos; el harness training los unifica. Y el costo de cómputo del run principal de CPT ---del orden de \textbf{USD 300--1,000} en GPU rentada--- es marginal frente al verdadero activo del proyecto: la ingeniería de datos sobre un corpus que ningún competidor posee.
\end{tesis}

El análisis de mercado muestra una oportunidad real y medible. El mercado LegalTech latinoamericano alcanzó USD 1.9 miles de millones en 2025 y se proyecta a USD 4.9 miles de millones en 2034 (CAGR 10.68\%),\cite{imarc2025} mientras el mercado global de servicios legales supera el billón de dólares anuales.\cite{mordorlegalservices2026} En México operan alrededor de 46,000 firmas de abogados\cite{industryresearch2026} y ejercen cerca de 449,000 abogados,\cite{mirada360} con más de 1.2 millones de egresados de derecho acumulados.\cite{forojuridico2023} La demanda ciudadana es masiva y está desatendida: la mitad de los mexicanos enfrentó un problema jurídico en los últimos dos años, pero solo 30\% buscó asesoría legal y apenas 5 de cada 100 acudieron a tribunales o policías.\cite{wjpmx2019,justiciaabierta2025} Estimamos un TAM mexicano de $\sim$USD 420M, un SAM de $\sim$USD 94M anuales en asientos profesionales, y un SOM a cinco años de USD 5--17M de ARR combinando suscripciones profesionales y despliegues soberanos on-premise.

La ventaja competitiva es defensible precisamente porque es incómoda de replicar: la regulación mexicana (LFPDPPP 2025, con la Secretaría Anticorrupción y Buen Gobierno como autoridad\cite{lfpdppp,ccn2025}) y el deber penal de secreto profesional (arts.~210--211 del CPF\cite{cpf}) hacen jurídicamente inaceptable para los despachos mexicanos subir documentos privilegiados a APIs extranjeras que pueden retener y reutilizar sus entradas. Los incumbentes globales ---Harvey (valorado en USD 11 mil millones\cite{harvey2026}), vLex Vincent, Thomson Reuters CoCounsel, Lexis+ AI--- operan exclusivamente en nube, con precios de USD 300--500 por asiento/mes,\cite{vaquill2026} y ninguno entrena sobre corpus mexicanos ni ofrece despliegue soberano. En este mercado, \textbf{el cumplimiento normativo no es una característica: es el foso defensivo}.

El documento se organiza como sigue: la Sección~\ref{sec:problema} define las tres fallas estructurales del estado del arte; la Sección~\ref{sec:necesidad} cuantifica la necesidad social y el mercado; la Sección~\ref{sec:marco} analiza el marco regulatorio que vuelve obligatoria la soberanía; la Sección~\ref{sec:relacionado} revisa el trabajo previo; las Secciones~\ref{sec:corpus} y~\ref{sec:entrenamiento} detallan el corpus y la estrategia de entrenamiento; la Sección~\ref{sec:viabilidad} presenta el modelo de negocio y sus unit economics; y las Secciones~\ref{sec:riesgos}--\ref{sec:roadmap} discuten riesgos, hoja de ruta y contribuciones.
